% Options for packages loaded elsewhere
% Options for packages loaded elsewhere
\PassOptionsToPackage{unicode}{hyperref}
\PassOptionsToPackage{hyphens}{url}
\PassOptionsToPackage{dvipsnames,svgnames,x11names}{xcolor}
%
\documentclass[
  letterpaper,
  DIV=11,
  numbers=noendperiod]{scrartcl}
\usepackage{xcolor}
\usepackage[margin=1in]{geometry}
\usepackage{amsmath,amssymb}
\setcounter{secnumdepth}{5}
\usepackage{iftex}
\ifPDFTeX
  \usepackage[T1]{fontenc}
  \usepackage[utf8]{inputenc}
  \usepackage{textcomp} % provide euro and other symbols
\else % if luatex or xetex
  \usepackage{unicode-math} % this also loads fontspec
  \defaultfontfeatures{Scale=MatchLowercase}
  \defaultfontfeatures[\rmfamily]{Ligatures=TeX,Scale=1}
\fi
\usepackage{lmodern}
\ifPDFTeX\else
  % xetex/luatex font selection
  \setmonofont[]{DejaVu Sans Mono}
\fi
% Use upquote if available, for straight quotes in verbatim environments
\IfFileExists{upquote.sty}{\usepackage{upquote}}{}
\IfFileExists{microtype.sty}{% use microtype if available
  \usepackage[]{microtype}
  \UseMicrotypeSet[protrusion]{basicmath} % disable protrusion for tt fonts
}{}
\makeatletter
\@ifundefined{KOMAClassName}{% if non-KOMA class
  \IfFileExists{parskip.sty}{%
    \usepackage{parskip}
  }{% else
    \setlength{\parindent}{0pt}
    \setlength{\parskip}{6pt plus 2pt minus 1pt}}
}{% if KOMA class
  \KOMAoptions{parskip=half}}
\makeatother
% Make \paragraph and \subparagraph free-standing
\makeatletter
\ifx\paragraph\undefined\else
  \let\oldparagraph\paragraph
  \renewcommand{\paragraph}{
    \@ifstar
      \xxxParagraphStar
      \xxxParagraphNoStar
  }
  \newcommand{\xxxParagraphStar}[1]{\oldparagraph*{#1}\mbox{}}
  \newcommand{\xxxParagraphNoStar}[1]{\oldparagraph{#1}\mbox{}}
\fi
\ifx\subparagraph\undefined\else
  \let\oldsubparagraph\subparagraph
  \renewcommand{\subparagraph}{
    \@ifstar
      \xxxSubParagraphStar
      \xxxSubParagraphNoStar
  }
  \newcommand{\xxxSubParagraphStar}[1]{\oldsubparagraph*{#1}\mbox{}}
  \newcommand{\xxxSubParagraphNoStar}[1]{\oldsubparagraph{#1}\mbox{}}
\fi
\makeatother

\usepackage{color}
\usepackage{fancyvrb}
\newcommand{\VerbBar}{|}
\newcommand{\VERB}{\Verb[commandchars=\\\{\}]}
\DefineVerbatimEnvironment{Highlighting}{Verbatim}{commandchars=\\\{\}}
% Add ',fontsize=\small' for more characters per line
\usepackage{framed}
\definecolor{shadecolor}{RGB}{241,243,245}
\newenvironment{Shaded}{\begin{snugshade}}{\end{snugshade}}
\newcommand{\AlertTok}[1]{\textcolor[rgb]{0.68,0.00,0.00}{#1}}
\newcommand{\AnnotationTok}[1]{\textcolor[rgb]{0.37,0.37,0.37}{#1}}
\newcommand{\AttributeTok}[1]{\textcolor[rgb]{0.40,0.45,0.13}{#1}}
\newcommand{\BaseNTok}[1]{\textcolor[rgb]{0.68,0.00,0.00}{#1}}
\newcommand{\BuiltInTok}[1]{\textcolor[rgb]{0.00,0.23,0.31}{#1}}
\newcommand{\CharTok}[1]{\textcolor[rgb]{0.13,0.47,0.30}{#1}}
\newcommand{\CommentTok}[1]{\textcolor[rgb]{0.37,0.37,0.37}{#1}}
\newcommand{\CommentVarTok}[1]{\textcolor[rgb]{0.37,0.37,0.37}{\textit{#1}}}
\newcommand{\ConstantTok}[1]{\textcolor[rgb]{0.56,0.35,0.01}{#1}}
\newcommand{\ControlFlowTok}[1]{\textcolor[rgb]{0.00,0.23,0.31}{\textbf{#1}}}
\newcommand{\DataTypeTok}[1]{\textcolor[rgb]{0.68,0.00,0.00}{#1}}
\newcommand{\DecValTok}[1]{\textcolor[rgb]{0.68,0.00,0.00}{#1}}
\newcommand{\DocumentationTok}[1]{\textcolor[rgb]{0.37,0.37,0.37}{\textit{#1}}}
\newcommand{\ErrorTok}[1]{\textcolor[rgb]{0.68,0.00,0.00}{#1}}
\newcommand{\ExtensionTok}[1]{\textcolor[rgb]{0.00,0.23,0.31}{#1}}
\newcommand{\FloatTok}[1]{\textcolor[rgb]{0.68,0.00,0.00}{#1}}
\newcommand{\FunctionTok}[1]{\textcolor[rgb]{0.28,0.35,0.67}{#1}}
\newcommand{\ImportTok}[1]{\textcolor[rgb]{0.00,0.46,0.62}{#1}}
\newcommand{\InformationTok}[1]{\textcolor[rgb]{0.37,0.37,0.37}{#1}}
\newcommand{\KeywordTok}[1]{\textcolor[rgb]{0.00,0.23,0.31}{\textbf{#1}}}
\newcommand{\NormalTok}[1]{\textcolor[rgb]{0.00,0.23,0.31}{#1}}
\newcommand{\OperatorTok}[1]{\textcolor[rgb]{0.37,0.37,0.37}{#1}}
\newcommand{\OtherTok}[1]{\textcolor[rgb]{0.00,0.23,0.31}{#1}}
\newcommand{\PreprocessorTok}[1]{\textcolor[rgb]{0.68,0.00,0.00}{#1}}
\newcommand{\RegionMarkerTok}[1]{\textcolor[rgb]{0.00,0.23,0.31}{#1}}
\newcommand{\SpecialCharTok}[1]{\textcolor[rgb]{0.37,0.37,0.37}{#1}}
\newcommand{\SpecialStringTok}[1]{\textcolor[rgb]{0.13,0.47,0.30}{#1}}
\newcommand{\StringTok}[1]{\textcolor[rgb]{0.13,0.47,0.30}{#1}}
\newcommand{\VariableTok}[1]{\textcolor[rgb]{0.07,0.07,0.07}{#1}}
\newcommand{\VerbatimStringTok}[1]{\textcolor[rgb]{0.13,0.47,0.30}{#1}}
\newcommand{\WarningTok}[1]{\textcolor[rgb]{0.37,0.37,0.37}{\textit{#1}}}

\usepackage{longtable,booktabs,array}
\newcounter{none} % for unnumbered tables
\usepackage{calc} % for calculating minipage widths
% Correct order of tables after \paragraph or \subparagraph
\usepackage{etoolbox}
\makeatletter
\patchcmd\longtable{\par}{\if@noskipsec\mbox{}\fi\par}{}{}
\makeatother
% Allow footnotes in longtable head/foot
\IfFileExists{footnotehyper.sty}{\usepackage{footnotehyper}}{\usepackage{footnote}}
\makesavenoteenv{longtable}
\usepackage{graphicx}
\makeatletter
\newsavebox\pandoc@box
\newcommand*\pandocbounded[1]{% scales image to fit in text height/width
  \sbox\pandoc@box{#1}%
  \Gscale@div\@tempa{\textheight}{\dimexpr\ht\pandoc@box+\dp\pandoc@box\relax}%
  \Gscale@div\@tempb{\linewidth}{\wd\pandoc@box}%
  \ifdim\@tempb\p@<\@tempa\p@\let\@tempa\@tempb\fi% select the smaller of both
  \ifdim\@tempa\p@<\p@\scalebox{\@tempa}{\usebox\pandoc@box}%
  \else\usebox{\pandoc@box}%
  \fi%
}
% Set default figure placement to htbp
\def\fps@figure{htbp}
\makeatother





\setlength{\emergencystretch}{3em} % prevent overfull lines

\providecommand{\tightlist}{%
  \setlength{\itemsep}{0pt}\setlength{\parskip}{0pt}}



 


\KOMAoption{captions}{tableheading}
\makeatletter
\@ifpackageloaded{caption}{}{\usepackage{caption}}
\AtBeginDocument{%
\ifdefined\contentsname
  \renewcommand*\contentsname{Table of contents}
\else
  \newcommand\contentsname{Table of contents}
\fi
\ifdefined\listfigurename
  \renewcommand*\listfigurename{List of Figures}
\else
  \newcommand\listfigurename{List of Figures}
\fi
\ifdefined\listtablename
  \renewcommand*\listtablename{List of Tables}
\else
  \newcommand\listtablename{List of Tables}
\fi
\ifdefined\figurename
  \renewcommand*\figurename{Figure}
\else
  \newcommand\figurename{Figure}
\fi
\ifdefined\tablename
  \renewcommand*\tablename{Table}
\else
  \newcommand\tablename{Table}
\fi
}
\@ifpackageloaded{float}{}{\usepackage{float}}
\floatstyle{ruled}
\@ifundefined{c@chapter}{\newfloat{codelisting}{h}{lop}}{\newfloat{codelisting}{h}{lop}[chapter]}
\floatname{codelisting}{Listing}
\newcommand*\listoflistings{\listof{codelisting}{List of Listings}}
\makeatother
\makeatletter
\makeatother
\makeatletter
\@ifpackageloaded{caption}{}{\usepackage{caption}}
\@ifpackageloaded{subcaption}{}{\usepackage{subcaption}}
\makeatother
\makeatletter
\@ifpackageloaded{tcolorbox}{}{\usepackage[skins,breakable]{tcolorbox}}
\makeatother
\makeatletter
\@ifundefined{shadecolor}{\definecolor{shadecolor}{rgb}{.97, .97, .97}}{}
\makeatother
\makeatletter
\makeatother
\makeatletter
\ifdefined\Shaded\renewenvironment{Shaded}{\begin{tcolorbox}[boxrule=0pt, breakable, enhanced, frame hidden, interior hidden, sharp corners]}{\end{tcolorbox}}\fi
\makeatother
\usepackage{bookmark}
\IfFileExists{xurl.sty}{\usepackage{xurl}}{} % add URL line breaks if available
\urlstyle{same}
\hypersetup{
  pdftitle={From Markdown to Blog Post: A ZZCOLLAB Conversion Workflow},
  pdfauthor={Ronald `Ryy' G. Thomas},
  colorlinks=true,
  linkcolor={blue},
  filecolor={Maroon},
  citecolor={Blue},
  urlcolor={Blue},
  pdfcreator={LaTeX via pandoc}}


\title{From Markdown to Blog Post: A ZZCOLLAB Conversion Workflow}
\usepackage{etoolbox}
\makeatletter
\providecommand{\subtitle}[1]{% add subtitle to \maketitle
  \apptocmd{\@title}{\par {\large #1 \par}}{}{}
}
\makeatother
\subtitle{Converting standalone documentation into reproducible blog
posts with Quarto}
\author{Ronald `Ryy' G. Thomas}
\date{2025-12-02}
\begin{document}
\maketitle

\renewcommand*\contentsname{Table of contents}
{
\setcounter{tocdepth}{3}
\tableofcontents
}

\begin{figure}[H]

{\centering \includegraphics[width=0.8\linewidth,height=\textheight,keepaspectratio]{media/images/hero.jpg}

}

\caption{A workspace with documentation ready for conversion into a
structured blog post}

\end{figure}%

\emph{Transforming standalone documentation into reproducible,
discoverable blog posts.}

\section{Introduction}\label{introduction}

Converting standalone markdown documentation into a professional blog
post presents a common challenge: a well-written technical reference
with working code examples and thorough coverage sits as a flat file on
disk, invisible to anyone who might benefit from it.

The gap between a markdown file and a published blog post turns out to
be narrower than expected, but only if a systematic workflow is
followed. The ZZCOLLAB framework provides the scaffolding: symlinks
bridge Quarto's expectations with the rrtools research compendium
layout, and a reproducible Docker environment ensures that readers can
rebuild the post from source.

This post documents that conversion workflow from start to finish,
including the decisions required, the directory structure to create, and
the verification steps that prevent broken links and missing images from
reaching the audience.

\subsection{Motivations}\label{motivations}

The following considerations motivated this workflow:

\begin{itemize}
\tightlist
\item
  Hundred-line markdown documents sitting in project directories remain
  invisible to anyone outside the project.
\item
  Search engines cannot index a markdown file on a local disk; a
  published blog post reaches a broader audience.
\item
  A repeatable process allows converting future documentation in minutes
  rather than hours.
\item
  The ZZCOLLAB dual-symlink system already solves the
  Quarto-versus-rrtools path conflict; documenting its use for
  documentation conversion specifically fills a gap.
\item
  A pre-publication checklist prevents publishing posts with broken
  image paths and missing YAML fields.
\end{itemize}

\subsection{Objectives}\label{objectives}

\begin{enumerate}
\def\labelenumi{\arabic{enumi}.}
\tightlist
\item
  Walk through the complete workflow for converting a standalone
  markdown file into a ZZCOLLAB-compatible Quarto blog post.
\item
  Explain the dual-symlink system that reconciles Quarto's directory
  expectations with the rrtools research compendium layout.
\item
  Provide concrete guidance on YAML frontmatter, media asset management,
  and content adaptation.
\item
  Deliver a pre-publication checklist covering content, metadata, media,
  structure, and testing.
\end{enumerate}

Errors and better approaches are welcome; see the Feedback section at
the end.

\begin{center}

\includegraphics[width=0.3\linewidth,height=\textheight,keepaspectratio]{media/images/quarto-logo.png}
\end{center}

\emph{A structured workflow turns documentation into discoverable
content.}

\section{Prerequisites and Setup}\label{prerequisites-and-setup}

You will need the following tools installed before starting:

\textbf{Required:}

\begin{itemize}
\tightlist
\item
  ZZCOLLAB framework (installed and accessible via \texttt{zzcollab}
  command)
\item
  Quarto (for rendering \texttt{.qmd} files to HTML)
\item
  Git (for version control)
\item
  Bash shell (macOS or Linux)
\end{itemize}

\textbf{Optional but useful:}

\begin{itemize}
\tightlist
\item
  GitHub CLI (\texttt{gh}) for automation
\item
  Docker (for full reproducibility)
\item
  A local Quarto preview server
\end{itemize}

\textbf{Background assumed:} Familiarity with markdown syntax, basic
command-line operations, and the concept of symlinks. No prior
experience with ZZCOLLAB is required; this post covers initialisation
from scratch.

\section{What is the ZZCOLLAB Blog
Workflow?}\label{what-is-the-zzcollab-blog-workflow}

The ZZCOLLAB blog workflow is a systematic process for converting
standalone markdown documentation into a Quarto blog post that lives
inside a reproducible research compendium. At its core, it solves a path
conflict: Quarto expects blog posts at \texttt{posts/*/index.qmd}, while
rrtools places authored content in \texttt{analysis/paper/}. The
dual-symlink system bridges both conventions so that a single source
file satisfies both tools.

Think of it as a translation layer. Your markdown content goes into the
rrtools-standard location. Symlinks at the project root make Quarto
believe the file is where it expects. Image paths like
\texttt{media/images/hero.jpg} resolve correctly from both locations
because symlinks exist in both directions.

The workflow itself takes roughly 40 minutes from start to finish: five
minutes for assessment, three for initialisation, one for symlink setup,
ten for content conversion, five for media assets, three for metadata,
five for testing, and two for committing.

\section{Getting Started}\label{getting-started}

\subsection{Step 1: Assess Your Source
Document}\label{step-1-assess-your-source-document}

Before touching any tools, spend five minutes evaluating the markdown
file you plan to convert. Ask yourself:

\begin{enumerate}
\def\labelenumi{\arabic{enumi}.}
\tightlist
\item
  How long is the document? A 200-line reference differs from a 600-line
  tutorial.
\item
  Does it contain code examples? If so, in which languages?
\item
  Does it need images? Hero images, diagrams, screenshots?
\item
  Is it a tutorial, reference guide, or how-to?
\item
  Who is the target audience?
\end{enumerate}

\begin{verbatim}
Example assessment:
- Length: ~600 lines (substantial)
- Code: Yes (bash scripts, R snippets)
- Images: Yes (hero + concept diagram)
- Type: How-to guide with reference sections
- Audience: R developers using ZZCOLLAB
- Estimate: 10 minutes to convert content
\end{verbatim}

\subsection{Step 2: Create the ZZCOLLAB
Project}\label{step-2-create-the-zzcollab-project}

\begin{Shaded}
\begin{Highlighting}[]
\BuiltInTok{cd}\NormalTok{ \textasciitilde{}/prj/qblog/posts}
\FunctionTok{mkdir}\NormalTok{ my\_blog\_post }\KeywordTok{\&\&} \BuiltInTok{cd}\NormalTok{ my\_blog\_post}
\ExtensionTok{zzcollab} \AttributeTok{{-}r}\NormalTok{ ubuntu\_standard\_publishing}
\end{Highlighting}
\end{Shaded}

This creates the standard project structure: \texttt{analysis/},
\texttt{R/}, \texttt{tests/}, a Dockerfile, \texttt{renv.lock}, a
Makefile, and \texttt{.zzcollab/manifest.json}.

\subsection{Step 3: Set Up the Dual-Symlink
System}\label{step-3-set-up-the-dual-symlink-system}

\begin{Shaded}
\begin{Highlighting}[]
\ExtensionTok{./modules/setup\_symlinks.sh}
\end{Highlighting}
\end{Shaded}

This creates symlinks at two levels:

\textbf{At the project root} (for Quarto discovery):

\begin{verbatim}
index.qmd  -> analysis/paper/index.qmd
figures/   -> analysis/figures/
media/     -> analysis/media/
data/      -> analysis/data/
\end{verbatim}

\textbf{Inside \texttt{analysis/paper/}} (for intuitive editing):

\begin{verbatim}
figures/ -> ../figures/
media/   -> ../media/
data/    -> ../data/
\end{verbatim}

The dual-symlink design means that when you write
\texttt{!{[}{]}(media/images/hero.jpg)} in your blog post, the path
resolves correctly whether Quarto renders from the project root or you
preview from \texttt{analysis/paper/}.

\begin{figure}[H]

{\centering \pandocbounded{
\includegraphics[keepaspectratio]{media/images/ambiance2.jpg}}

}

\caption{An organised bookshelf with reference materials and a notebook}

\end{figure}%

\emph{Proper directory structure prevents the most common blog post
failures.}

\section{Converting Your Content}\label{converting-your-content}

\subsection{YAML Frontmatter}\label{yaml-frontmatter}

Replace the template YAML with metadata specific to your post. The
fields that matter most for blog listing display are \texttt{title},
\texttt{categories}, \texttt{description}, and \texttt{image}:

\begin{Shaded}
\begin{Highlighting}[]
\PreprocessorTok{{-}{-}{-}}
\FunctionTok{title}\KeywordTok{:}\AttributeTok{ }\StringTok{"Your Blog Post Title"}
\FunctionTok{subtitle}\KeywordTok{:}\AttributeTok{ }\StringTok{"Clarifying subtitle"}
\FunctionTok{author}\KeywordTok{:}\AttributeTok{ }\StringTok{"Your Name"}
\FunctionTok{date}\KeywordTok{:}\AttributeTok{ }\StringTok{"2025{-}12{-}02"}
\FunctionTok{categories}\KeywordTok{:}\AttributeTok{ }\KeywordTok{[}\AttributeTok{quarto}\KeywordTok{,}\AttributeTok{ zzcollab}\KeywordTok{,}\AttributeTok{ workflow}\KeywordTok{,}\AttributeTok{ reproducibility}\KeywordTok{]}
\FunctionTok{description}\KeywordTok{:}\AttributeTok{ }\StringTok{"2{-}3 sentence summary."}
\FunctionTok{image}\KeywordTok{:}\AttributeTok{ }\StringTok{"media/images/hero.jpg"}
\FunctionTok{document{-}type}\KeywordTok{:}\AttributeTok{ }\StringTok{"blog"}
\FunctionTok{draft}\KeywordTok{:}\AttributeTok{ }\CharTok{false}
\FunctionTok{execute}\KeywordTok{:}
\AttributeTok{  }\FunctionTok{echo}\KeywordTok{:}\AttributeTok{ }\CharTok{true}
\AttributeTok{  }\FunctionTok{warning}\KeywordTok{:}\AttributeTok{ }\CharTok{false}
\AttributeTok{  }\FunctionTok{message}\KeywordTok{:}\AttributeTok{ }\CharTok{false}
\FunctionTok{format}\KeywordTok{:}
\AttributeTok{  }\FunctionTok{html}\KeywordTok{:}
\AttributeTok{    }\FunctionTok{code{-}fold}\KeywordTok{:}\AttributeTok{ }\CharTok{false}
\AttributeTok{    }\FunctionTok{code{-}tools}\KeywordTok{:}\AttributeTok{ }\CharTok{false}
\PreprocessorTok{{-}{-}{-}}
\end{Highlighting}
\end{Shaded}

Key points:

\begin{itemize}
\tightlist
\item
  Keep titles under 70 characters and include searchable keywords.
\item
  Use 2-4 categories drawn from your blog's existing taxonomy.
\item
  The \texttt{description} field appears in blog listings; lead with the
  value proposition.
\item
  Set \texttt{draft:\ false} during development.
\end{itemize}

\subsection{Adapting Markdown to
Quarto}\label{adapting-markdown-to-quarto}

Most markdown syntax transfers directly. The main conversions to
consider:

\begin{itemize}
\tightlist
\item
  \textbf{Blockquotes} become Quarto callouts:
\end{itemize}

\begin{Shaded}
\begin{Highlighting}[]
\NormalTok{::: \{.callout{-}note\}}
\NormalTok{\#\# Summary}
\NormalTok{Your key takeaway goes here.}
\NormalTok{:::}
\end{Highlighting}
\end{Shaded}

\begin{itemize}
\tightlist
\item
  \textbf{Code blocks} should specify a language for syntax highlighting
  (\texttt{bash}, \texttt{r}, \texttt{python}).
\item
  \textbf{Images} use the \texttt{.img-fluid} class for responsive
  display.
\end{itemize}

\subsection{Adding Media Assets}\label{adding-media-assets}

Copy images into the media directory and document their sources:

\begin{Shaded}
\begin{Highlighting}[]
\FunctionTok{cp}\NormalTok{ \textasciitilde{}/path/to/hero.jpg }\DataTypeTok{\textbackslash{}}
\NormalTok{  analysis/media/images/hero.jpg}

\FunctionTok{cat} \OperatorTok{\textgreater{}}\NormalTok{ analysis/media/images/README.md }\OperatorTok{\textless{}\textless{} \textquotesingle{}EOF\textquotesingle{}}
\StringTok{\# Image Attribution}

\StringTok{\#\# hero.jpg}
\StringTok{{-} Source: Unsplash}
\StringTok{{-} Photographer: [Name]}
\StringTok{{-} Licence: Unsplash Licence}
\StringTok{{-} URL: [link]}
\OperatorTok{EOF}
\end{Highlighting}
\end{Shaded}

Reference images in your post with relative paths:

\begin{Shaded}
\begin{Highlighting}[]
\NormalTok{![Alt text](media/images/hero.jpg)\{.img{-}fluid\}}

\NormalTok{*Caption with attribution.*}
\end{Highlighting}
\end{Shaded}

\section{Integration with the Parent
Blog}\label{integration-with-the-parent-blog}

The parent Quarto blog configuration (\texttt{\_quarto.yml}) already
renders posts matching the pattern \texttt{posts/*/index.qmd}. Because
the root-level symlink maps \texttt{index.qmd} to
\texttt{analysis/paper/index.qmd}, the blog discovers your post
automatically.

No changes to the parent blog configuration are required. When you
render the entire site with \texttt{quarto\ render}, the new post
appears in the listing.

\section{Version Control}\label{version-control}

Commit your work with a descriptive message that documents the content,
media assets, and reproducibility infrastructure:

\begin{Shaded}
\begin{Highlighting}[]
\BuiltInTok{cd}\NormalTok{ \textasciitilde{}/prj/qblog/posts/my\_blog\_post}
\FunctionTok{git}\NormalTok{ add .}
\FunctionTok{git}\NormalTok{ status}
\FunctionTok{git}\NormalTok{ commit }\AttributeTok{{-}m} \StringTok{"Add blog post: Your Post Title}

\StringTok{{-} analysis/paper/index.qmd: Main content}
\StringTok{{-} analysis/media/images/: Hero and diagrams}
\StringTok{{-} Dockerfile + renv.lock: Reproducible env"}
\end{Highlighting}
\end{Shaded}

Git handles symlinks correctly; they are stored as text files containing
the target path.

\section{Testing and Validation}\label{testing-and-validation}

\subsection{Quick Validation}\label{quick-validation}

\begin{Shaded}
\begin{Highlighting}[]
\FunctionTok{ls} \AttributeTok{{-}la}\NormalTok{ index.qmd}
\FunctionTok{head} \AttributeTok{{-}25}\NormalTok{ analysis/paper/index.qmd}
\FunctionTok{grep} \StringTok{"images/"}\NormalTok{ analysis/paper/index.qmd}
\end{Highlighting}
\end{Shaded}

\subsection{Local Render}\label{local-render}

\begin{Shaded}
\begin{Highlighting}[]
\BuiltInTok{cd}\NormalTok{ \textasciitilde{}/prj/qblog/posts/my\_blog\_post}
\ExtensionTok{quarto}\NormalTok{ render analysis/paper/index.qmd}
\ExtensionTok{open}\NormalTok{ index.html}
\end{Highlighting}
\end{Shaded}

Common issues and fixes:

{\def\LTcaptype{none} % do not increment counter
\begin{longtable}[]{@{}lll@{}}
\toprule\noalign{}
Error & Cause & Fix \\
\midrule\noalign{}
\endhead
\bottomrule\noalign{}
\endlastfoot
\texttt{index.qmd\ not\ found} & Broken symlink &
\texttt{ls\ -la\ index.qmd} \\
Image not found & Wrong path & Use \texttt{media/images/file.jpg} \\
YAML parse error & Tabs in YAML & Use spaces only \\
No highlighting & Missing language & Specify
\texttt{\textasciigrave{}\textasciigrave{}\textasciigrave{}bash} \\
\end{longtable}
}

\subsection{Things to Watch Out For}\label{things-to-watch-out-for}

\begin{enumerate}
\def\labelenumi{\arabic{enumi}.}
\tightlist
\item
  \textbf{Symlinks break after cloning.} If someone clones your
  repository and the symlinks do not recreate correctly, they must run
  \texttt{setup\_symlinks.sh} again or recreate them manually.
\item
  \textbf{Tabs in YAML cause silent failures.} Quarto YAML requires
  spaces for indentation. A single tab character produces a parse error
  that does not always point to the offending line.
\item
  \textbf{Absolute image paths break portability.} Always use relative
  paths (\texttt{media/images/hero.jpg}), never absolute paths. Absolute
  paths work on your machine and nowhere else.
\item
  \textbf{The \texttt{draft} field defaults to \texttt{false} when
  omitted.} If you forget to set \texttt{draft:\ false}, your unfinished
  post appears on the public site after the next render.
\item
  \textbf{Large images slow rendering and page load.} Keep images under
  500 KB each. Use JPEG for photographs and PNG for diagrams with sharp
  edges.
\end{enumerate}

\begin{figure}[H]

{\centering \pandocbounded{\includegraphics[keepaspectratio]{media/images/ambiance3.jpg}}

}

\caption{A notebook with a completed checklist beside a cup of coffee}

\end{figure}%

\emph{A pre-publication checklist prevents common deployment failures.}

\subsection{Lessons Learnt}\label{lessons-learnt}

\subsubsection{Conceptual Understanding}\label{conceptual-understanding}

\begin{itemize}
\tightlist
\item
  The dual-symlink system is the single most important architectural
  decision in the workflow; it reconciles two competing directory
  conventions without requiring changes to either tool.
\item
  Blog posts are more discoverable than standalone markdown files; the
  effort of conversion pays off in audience reach.
\item
  Structured YAML metadata enables Quarto's listing features (search,
  category filtering, date sorting) without custom JavaScript.
\item
  The rrtools research compendium layout provides a natural separation
  between authored content, analysis code, and media assets.
\end{itemize}

\subsubsection{Technical Skills}\label{technical-skills}

\begin{itemize}
\tightlist
\item
  Setting up symlinks with \texttt{ln\ -s} and verifying them with
  \texttt{ls\ -la} became second nature after the first conversion.
\item
  Writing Quarto callouts (\texttt{:::\ \{.callout-note\}\ ...\ :::})
  replaced blockquotes for emphasis and improved visual hierarchy.
\item
  The \texttt{.img-fluid} class in Bootstrap ensures images scale
  correctly on mobile without any custom CSS.
\item
  Using \texttt{quarto\ render} for local preview before committing
  caught broken paths that would otherwise reach the public site.
\end{itemize}

\subsubsection{Gotchas and Pitfalls}\label{gotchas-and-pitfalls}

\begin{itemize}
\tightlist
\item
  Forgetting to document image sources in
  \texttt{analysis/media/images/README.md} requires reconstructing
  attribution after the fact.
\item
  A missing \texttt{.img-fluid} class causes hero images to overflow
  their container on mobile devices.
\item
  Committing a post with \texttt{draft:\ false} by accident results in
  the post appearing live on the site before the metadata is corrected.
\item
  Quarto's error messages for broken image paths are sometimes vague;
  using \texttt{grep} for all image references before rendering saves
  debugging time.
\end{itemize}

\subsection{Limitations}\label{limitations}

\begin{itemize}
\tightlist
\item
  This workflow assumes a Unix-like environment (macOS or Linux) for
  symlink support. Windows users may need WSL or alternative approaches.
\item
  The 40-minute estimate applies to straightforward documentation with
  minimal code and few images. Complex tutorials with interactive
  elements take longer.
\item
  The dual-symlink system adds conceptual complexity that may confuse
  contributors unfamiliar with the convention.
\item
  ZZCOLLAB must be installed and functional before starting.
  Troubleshooting ZZCOLLAB itself is outside the scope of this workflow.
\item
  The workflow does not cover adding R analysis pipelines to posts that
  are purely narrative; extending a documentation post with data
  analysis requires additional steps.
\item
  Local rendering requires Quarto installed on the author's machine.
  Docker-based rendering is possible but adds setup time.
\end{itemize}

\subsection{Opportunities for
Improvement}\label{opportunities-for-improvement}

\begin{enumerate}
\def\labelenumi{\arabic{enumi}.}
\tightlist
\item
  Create a shell script that automates the entire workflow: directory
  creation, ZZCOLLAB initialisation, symlink setup, and YAML template
  population from command-line arguments.
\item
  Build a validation script that checks all image paths, verifies
  symlink integrity, and confirms YAML field completeness before
  rendering.
\item
  Add a \texttt{make\ convert} target to the Makefile that accepts a
  source markdown path and performs the conversion automatically.
\item
  Develop a Quarto extension that validates \texttt{document-type}
  metadata and warns when required sections are missing.
\item
  Create a companion post covering the reverse workflow: extracting
  documentation from an existing blog post into a standalone reference.
\item
  Implement automated image optimisation (resizing, compression) as part
  of the media asset pipeline.
\end{enumerate}

\section{Wrapping Up}\label{wrapping-up}

Converting standalone markdown documentation into a ZZCOLLAB blog post
is a systematic, repeatable process that takes roughly 40 minutes once
you understand the workflow. The dual-symlink system bridges the gap
between Quarto's directory expectations and the rrtools research
compendium layout, and proper YAML metadata ensures your content is
discoverable through search and category filtering.

The key insight from developing this workflow is that the barrier to
publishing is not the conversion itself but the lack of a structured
process. Once the steps are documented and the checklist is in place,
each subsequent conversion is faster and less error-prone.

For those with documentation sitting in project directories that could
benefit a broader audience, this workflow merits consideration. The
initial investment in understanding the symlink system pays dividends
every time another document is converted.

Main takeaways:

\begin{itemize}
\tightlist
\item
  The dual-symlink system is the architectural foundation; understand it
  first.
\item
  YAML metadata drives discoverability; invest time in titles,
  descriptions, and categories.
\item
  A pre-publication checklist prevents the most common failures.
\item
  The entire process takes about 40 minutes from assessment to committed
  post.
\end{itemize}

\section{See Also}\label{see-also}

\textbf{Related posts:}

\begin{itemize}
\tightlist
\item
  \href{../39-templatepost/templatepost/docs/\%20ZZCOLLAB_BLOG_SETUP.md}{ZZCOLLAB
  Blog Setup Guide}: Comprehensive setup documentation
\item
  \href{../39-templatepost/templatepost/analysis/report/\%20index.qmd}{Blog
  Post Template}: The template this workflow produces
\end{itemize}

\textbf{Key resources:}

\begin{itemize}
\tightlist
\item
  \href{https://github.com/mygit/zzcollab}{ZZCOLLAB Framework}: Project
  initialisation and management
\item
  \href{https://quarto.org/docs/websites/website-blog.html}{Quarto Blog
  Guide}: Official Quarto blog documentation
\item
  \href{https://github.com/benmarwick/rrtools}{rrtools Research
  Compendium}: The convention ZZCOLLAB follows
\item
  \href{https://quarto.org/docs/reference/formats/html.html}{Quarto YAML
  Reference}: Complete YAML options
\end{itemize}

\section{Reproducibility}\label{reproducibility}

This post is primarily narrative and does not involve a computational
analysis pipeline. The ZZCOLLAB project structure provides the
reproducibility infrastructure:

\begin{Shaded}
\begin{Highlighting}[]
\BuiltInTok{cd}\NormalTok{ \textasciitilde{}/prj/qblog/posts/06{-}markdowntoblog/}
\BuiltInTok{cd}\NormalTok{ markdowntoblog}
\ExtensionTok{quarto}\NormalTok{ render analysis/report/index.qmd}
\end{Highlighting}
\end{Shaded}

\textbf{Project files:}

\begin{verbatim}
markdowntoblog/
  analysis/report/index.qmd   # This post
  analysis/media/images/      # Hero, ambiance
  Dockerfile                  # Environment
  renv.lock                   # R packages
  Makefile                    # Build automation
\end{verbatim}

\section{Feedback}\label{feedback}

Feedback is welcome for:

\begin{itemize}
\tightlist
\item
  Errors or better approaches to any of the steps in this workflow
\item
  Suggestions for topics to cover in future posts
\item
  Discussion of documentation workflows, reproducibility, or Quarto
  publishing
\item
  Questions about ZZCOLLAB or the symlink system
\item
  General discussion of technical writing for data science
\end{itemize}




\end{document}
